\section{来源审计与课程主问题}

这讲课的中心问题不是“参数量是否重要”，而是：当一个语言模型能够流畅生成答案时，我们还需要哪些机制，才能把它变成更可靠的推理器、知识模型或伦理判断组件？Yejin Choi 用 Maieutic Prompting、Symbolic Knowledge Distillation 和 Delphi 三项研究回答这个问题。三者共同依赖大型语言模型，却都在模型外增加结构、数据、critic 或约束，因此不能被压缩成一句“小模型打败大模型”。

\subsection{本讲如何重建}

课程官网列出录像与三篇指定论文，但没有公开独立 slide deck。本讲以 Stanford Online 1080p 官方录像为视觉母版：每两秒采样得到 2,253 帧，稳定帧工具产生 207 个高召回候选，再人工去重为 56 个独立教学状态。最终图像均从 1920x1080 全宽画面提取；52:40 之后的问答大量回翻旧页，因此重复图不再插入，口头解释则通过官方 \texttt{en-US} 人工字幕保留。

\begin{longtable}{@{}L{0.18\textwidth}L{0.27\textwidth}L{0.43\textwidth}@{}}
\toprule
证据层 & 本地材料 & 在讲义中的职责 \\
\midrule
官方录像 & 1:15:05 视频与 56 张 teaching states & 决定课堂顺序、公式页、结果表、公开演示、争议回应与 hybrid 架构 \\
官方人工字幕 & \texttt{lecture16.en.srt} 与 timed transcript & 恢复讲者的动机、比较 caveat、数据偏差、系统成本和 Q\&A \\
指定论文 & Maieutic、Symbolic KD、Delphi & 核对算法、公式、实验范围、数据构成与限制 \\
背景论文 & COMET、COMET-ATOMIC 2020 & 解释 symbolic graph 与 neural knowledge model 的关系 \\
审计文件 & manifest、coverage、ledger、blueprint & 证明视觉和 teacher voice 节点都有明确教学去向 \\
\bottomrule
\end{longtable}

\begin{warningbox}{历史边界：2023 年 2 月 14 日}
本讲讨论的是 ChatGPT 刚公开数月时的系统与争论。后来的 reasoning models、post-training 方法和 machine ethics 产品不能倒灌为课堂证据。尤其是 Delphi hybrid 在课上仍是 work in progress，不能写成已经成熟部署的系统。
\end{warningbox}

\subsection{David 与 Goliath：标题不是反 Scaling 宣言}

标题页把极大模型比作 Goliath，把更小但结构化的系统比作 David。读这张图时要先保留限定词：讲者承认大模型在许多任务上更强，只是反对从少数流畅例子推导出“常识已经解决”。本讲要找的是另一条增益路径：更真实的 evaluation、更针对性的 task、更高质量的数据，以及能检查语言模型输出的算法。

\lecturefigure{slide-01-title-david-vs-goliath.jpg}{David vs. Goliath：极大语言模型时代的排行榜艺术}{官方录像恢复幻灯片 V001；Stanford CS25 Lecture 16}

下一张 slide 借《孙子兵法》把研究方法拆成三步：先认识失败分布，再选择真正有区分度的战场，最后设计新的武器。它不是文学装饰，而是整堂课的 evidence discipline：如果 benchmark 只奖励 dataset artifact，再高的分数也不能证明模型掌握了 underlying task。

\lecturefigure{slide-02-art-of-war-lessons.jpg}{Know your enemy、choose your battles、innovate your weapons}{官方录像恢复幻灯片 V002；Stanford CS25 Lecture 16}

\teachervoice{讲者用“existential crisis”自嘲 ChatGPT 带来的冲击，但随即把问题改写为 hasty generalization：一次答对 Winograd 题不等于拥有稳定物理常识。她反复强调的不是“规模无用”，而是“smaller can be better”与“knowledge is power”。}

\subsection{术语表：先把相近术语分开}

后文会密集出现 language model、knowledge model、commonsense graph、critic、descriptive morality 与 normative ethics。若不先区分对象，读者很容易把“生成一句合理文本”“拥有可复用知识”“做出伦理判断”混成同一种能力。下面的表给出本讲使用的工作定义。

\begin{longtable}{@{}L{0.19\textwidth}L{0.27\textwidth}L{0.42\textwidth}@{}}
\toprule
术语 & 解决的问题 & 核心机制与课程关系 \\
\midrule
Language model & 下一个 token 如何分布 & 从文本统计学习 $p(x_t\mid x_{<t})$；流畅性与世界知识可能相关，但不等价 \\
Knowledge model & 某个实体、事件或关系能推出什么 & 以 typed relation、检索接口或专门生成目标保存可复用知识 \\
Commonsense & 日常情境中广泛共享的实用知识与推理 & 包含物理、社会、因果、意图与反应；并非所有文化都完全一致 \\
Symbolic graph & 如何显式表示节点和关系 & ATOMIC 用事件节点与 relation type 形成可检查的知识结构 \\
Critic & 哪些机器生成内容应被保留 & 学习质量判别并过滤 teacher samples；critic 本身仍会犯错 \\
Descriptive ethics & 人们通常会怎样判断 & Delphi 预测 sampled human judgments，不宣称给出终极正确答案 \\
Normative ethics & 人们应该怎样判断 & 需要原则、论证与制度参与，不能由多数标签自动推出 \\
Reflective equilibrium & 如何协调案例直觉与抽象原则 & 在 bottom-up judgment 与 top-down constraint 之间反复修正 \\
Adversarial test & 系统在反常、组合或蓄意构造输入下是否稳定 & 用来暴露 dataset success 与 task understanding 的差距 \\
\bottomrule
\end{longtable}

\subsection{四个贯穿全讲的问题}

阅读后续三套系统时，应持续追问四件事。第一，base LM 提供的是 fluent generation、latent knowledge 还是 calibrated truth？第二，系统把什么结构放到模型外：logic、knowledge graph、critic 还是 moral principle？第三，结果究竟测 accuracy、consistency、human plausibility 还是 deployment safety？第四，失败能否被定位并修复，还是只得到另一个不可解释的总分？

\begin{importantbox}{本讲的共同 Recipe}
先让大模型产生候选，再用显式结构或专用模型缩小候选空间，最后把结果训练进更小、更专门的系统。Maieutic 用 logic 约束 explanations，Symbolic KD 用 critic 约束 knowledge corpus，Delphi / hybrid 用 commonsense data 与 moral constraints 约束 judgment。
\end{importantbox}

\subsection{本章小结}

来源审计表明，这是一讲关于系统分工与证据边界的课程。David-versus-Goliath 不是否定规模，而是要求把语言模型放进可检查的 pipeline；后续每一章都要说明候选从哪里来、约束由谁施加、指标证明了什么，以及指标没有证明什么。

